% !Mode:: "TeX:UTF-8"
% 物理量名称及符号表。一行一个符号：\item[符号] 说明。
%
% 第一列宽度按内容定（2.5cm 起，挑整张表最矮的那个宽）；要钉死就写 \begin{symbols}[3cm]，
% 也可以写成键值 \begin{symbols}[labelwidth=3cm]。
%
% 规范里允许把符号分成几张表、各有各的列。列表排不出那种版式，要的话把环境名加个
% 星号，正文自己写表格：
%
%   \begin{symbols*}
%   \begin{longtblr}{colspec={cccc}} ... \end{longtblr}
%   \end{symbols*}

\begin{symbols}
\item[$p$]              气膜压力，Pa
\item[$h$]              气膜厚度，m
\item[$\eta$]           气体动力黏度，Pa$\cdot$s
\item[$D_{\mathrm{p}}$] 多孔质材料的平均粒子直径，m
\item[$\psi$]           孔隙度
\item[$C_2$]            特征渗透性或固有渗透性，m$^2$
\end{symbols}
